\chapter{Dynamics: Physical Necessity}

\otbookepigraph{Dynamics is the study of why things move, not just where they go.}{OT Interpretation}

\section{De-construction: Euler-Lagrange}
The \acr{el} equation is:
\begin{equation}
    \frac{d}{dt}\frac{\partial L}{\partial \dot{q}} - \frac{\partial L}{\partial q} = \tau
\end{equation}

\section{Re-interpretation: The ``Budget'' Analogy}
\begin{intuitionbox}[The Energy Budget]
    Nature is lazy. It wants to minimize the ``Action'' (the integral of \(L = T-V\)).
    Think of \(T\) (Kinetic) as your current cash and \(V\) (Potential) as your debt. The \acr{el} equation is the bookkeeping rule that tells you how much force \(\tau\) you need to apply to move between states while respecting your energy budget.
\end{intuitionbox}

\section{The Inertia Matrix \sym{mass-matrix}}
\begin{designrule}[Matrix Invariants]
    The inertia matrix \sym{mass-matrix} is always symmetric and positive-definite.
    \textit{Interpretation:} A robot cannot have negative mass, and if you push it, it must move in a way that consumes energy. If your derived \sym{mass-matrix} is not positive-definite, your model is physically impossible.
\end{designrule}

\section{Mental Re-construction}
\begin{practice}[The 1-Link Test]
    Derive the \acr{el} equation for a single pendulum. Identify the \sym{gravity-g} term. Why is it a sine function? (Because gravity acts on the lever arm's vertical displacement).
\end{practice}
